\chapter*{Introduction Générale}
\addcontentsline{toc}{chapter}{Introduction Générale}

Dans un contexte économique mondial marqué par une transformation numérique
accélérée, le Système d'Information (SI) est devenu une composante stratégique
incontournable de toute organisation. Comme le souligne Reix~\cite{reix1995}, « un
système d'information peut être défini comme un ensemble organisé de ressources
(matériels, logiciels, personnel et procédures) permettant d'acquérir, de traiter, de stocker
et de communiquer les informations dans les organisations ». À ce titre, sa disponibilité
et sa fiabilité conditionnent directement la performance opérationnelle de l'entreprise.

Naftal, filiale à 100\% de Sonatrach et leader national de la commercialisation et
de la distribution des produits pétroliers en Algérie, opère l'un des SI industriels les
plus étendus du pays. Ce SI supporte des activités critiques : gestion des stocks de
carburants, suivi d'une flotte de 3~600 véhicules, coordination de 1~846 stations-service
et administration d'un réseau intranet reliant l'ensemble des branches et districts sur
tout le territoire national. La continuité de ce SI est donc une exigence absolue.

Or, la gestion des pannes et incidents informatiques au sein du Groupe Informatique
de la Branche Carburants s'effectue encore aujourd'hui de manière traditionnelle : appels
téléphoniques, courriers électroniques non structurés, absence de base de connaissances.
Cette approche manuelle engendre des pertes d'informations, une désorganisation des
équipes techniques et une impossibilité de mesurer objectivement la qualité du support.

C'est dans ce contexte précis que s'inscrit notre projet de fin d'études.

\medskip
\noindent\textbf{La problématique générale peut se formuler ainsi :} \emph{Comment doter le Groupe
Informatique de la Branche Carburants de Naftal d'un outil de support applicatif centralisé,
conforme aux contraintes d'infrastructure de l'entreprise, et permettant une traçabilité complète
ainsi qu'un pilotage objectif de la qualité de service ?}
\medskip

Cette question prolonge, à l'échelle d'une équipe de proximité, les préoccupations identifiées
dans les travaux sur la performance de la fonction SI et de la DCSI~\cite{meghari2014} : lorsque
l'assistance reste informelle, la fonction support peine à démontrer sa valeur et à s'améliorer.

L'objectif opérationnel est de concevoir et de réaliser une \textbf{Plateforme Intégrée de Gestion
et de Suivi des Incidents Applicatifs (GIA)}, solution de type \textit{Helpdesk} web intranet
permettant de centraliser, tracer et piloter le cycle de vie des incidents applicatifs concernés.

Ce mémoire s'articule autour des volets suivants :
\begin{itemize}
    \item \textbf{Chapitre 1 — Cadre théorique :} Rappels sur le système d'information,
    ses fonctions et la place de la fonction support dans la grande entreprise.
    \item \textbf{Chapitre 2 — Contexte et existant :} Présentation de Naftal, du SI et
    de la Branche Carburants ; critique du processus actuel de gestion des incidents.
    \item \textbf{Chapitre 3 — Analyse des besoins :} Acteurs, exigences fonctionnelles et
    non fonctionnelles, modélisation UML (cas d'utilisation et compléments).
    \item \textbf{Chapitre 4 — Conception :} Architecture 3-tiers, cycle de vie des tickets,
    modèle de données.
    \item \textbf{Chapitre 5 — Réalisation :} Environnement Windows Server / IIS, choix
    technologiques et interfaces de la plateforme GIA.
\end{itemize}
Une \textbf{conclusion générale} synthétise les apports du projet et propose des perspectives
d'évolution.
